\hojaejercicio{Parte 1: Capítulo 2}
\setcounter{ejercicio}{0}

% Ejercicio 1
\ejercicio{resuelto}{}{
    Suponiendo que una fuerza friccional resiste el movimiento del péndulo simple y que esa fuerza es proporcional, con un coeficiente $\kappa > 0$, a la velocidad del péndulo, la ecuación del péndulo simple se convierte en
    \[
        MLu'' + \kappa Lu' + MG \sin u = 0.
    \]
    Analizar la existencia, unicidad y comportamiento global de las soluciones maximales del problema de Cauchy asociado a esta ecuación diferencial.
}{
    Dividiendo la ecuación por $ML > 0$ y definiendo el coeficiente de amortiguamiento $r = \frac{\kappa}{M} > 0$ y la frecuencia propia al cuadrado $\omega_0^2 = \frac{G}{L} > 0$, el problema de Cauchy general con datos iniciales en $t = t_0$ adopta la forma canónica:
    \[
        (P) \quad
        \begin{cases}
            u'' = - r u' - \frac{G}{L} \sin u, \\
            u(t_0) = u_0, \quad u'(t_0) = v_0,
        \end{cases}
        \qquad r = \frac{\kappa}{M} > 0, \quad \frac{G}{L} > 0.
    \]

    \textbf{1. Reducción a sistema de primer orden y regularidad:}

    Introduciendo las variables de estado en el plano de fases $(u, v) = (u, u')$, la ecuación escalar de segundo orden equivale al sistema autónomo de primer orden en $\mathbb{R}^2$:
    \[
        \begin{cases}
            u' = v, \\
            v' = - r v - \frac{G}{L} \sin u,
        \end{cases}
        \iff
        \begin{pmatrix} u \\ v \end{pmatrix}' = \begin{pmatrix} v \\ - r v - \frac{G}{L} \sin u \end{pmatrix} = f(u, v),
    \]
    con dato inicial $(u(t_0), v(t_0)) = (u_0, v_0) \in \mathbb{R}^2$.
    
    El campo vectorial $f: \mathbb{R}^2 \to \mathbb{R}^2$ no depende explícitamente de $t$ (autónomo) y sus componentes son funciones analíticas ($\mathcal{C}^\infty$). Su matriz jacobiana respecto de $(u, v)$ es:
    \[
        D_{(u,v)} f(u, v) = \begin{pmatrix}
            0 & 1 \\
            - \frac{G}{L} \cos u & -r
        \end{pmatrix}.
    \]

    \textbf{2. Carácter Lipschitziano global y existencia y unicidad:}

    Analizamos la norma de las entradas de la matriz jacobiana en todo el espacio de fases $\mathbb{R}^2$:
    \[
        |0| = 0, \quad |1| = 1, \quad |-r| = r, \quad \left| - \frac{G}{L} \cos u \right| \le \frac{G}{L} \qquad \forall (u, v) \in \mathbb{R}^2.
    \]
    Al estar la matriz $D_{(u,v)} f(u, v)$ uniformemente acotada en todo $\mathbb{R}^2$, el campo vectorial $f$ es \textbf{globalmente Lipschitziano} en $\mathbb{R}^2$, con constante de Lipschitz independiente de $t$.

    Por el Teorema de Cauchy-Lipschitz global (o considerando una sucesión de intervalos expansivos $J_n = [t_0 - n, t_0 + n]$ y pasando al límite gracias a la unicidad), para cualquier dato inicial $(u_0, v_0) \in \mathbb{R}^2$ existe una \textbf{única solución maximal} $(u, v) \in \mathcal{C}^1(\mathbb{R}; \, \mathbb{R}^2)$, globalmente definida en todo tiempo:
    \[
        I = (-\infty, +\infty).
    \]
    Por tanto, $u \in \mathcal{C}^2(\mathbb{R}; \, \mathbb{R})$ existe de forma global sin posibilidad de explosión en tiempo finito ($T_{\max} = +\infty$, $T_{\min} = -\infty$).

    \medskip

    \textbf{3. Comportamiento global hacia el futuro (disipación de energía):}

    Definimos la energía mecánica total del péndulo (cinética más potencial gravitatoria):
    \[
        E(u, v) = \frac{1}{2} v^2 + \frac{G}{L} (1 - \cos u) \ge 0 \qquad \forall (u, v) \in \mathbb{R}^2.
    \]
    Derivando la energía a lo largo de cualquier trayectoria solución $(u(t), v(t))$:
    \[
        \frac{d}{dt} E(u(t), v(t)) = v(t) v'(t) + \frac{G}{L} \sin u(t) u'(t) = v(t) \left( -r v(t) - \frac{G}{L} \sin u(t) \right) + \frac{G}{L} v(t) \sin u(t) = - r v^2(t) \le 0.
    \]
    Puesto que $r = \frac{\kappa}{M} > 0$, la energía mecánica es una función monótona no creciente del tiempo. En particular:
    \[
        \frac{1}{2} v^2(t) \le E(u(t), v(t)) \le E(u_0, v_0) \qquad \forall t \ge t_0,
    \]
    lo que garantiza que la velocidad angular permanece uniformemente acotada para todo $t \ge t_0$:
    \[
        |u'(t)| = |v(t)| \le \sqrt{2 E(u_0, v_0)} \qquad \forall t \ge t_0.
    \]
    Al ser un sistema disipativo con rozamiento estricto, la energía decrece estrictamente siempre que $v(t) \neq 0$. Por el Principio de Invarianza de LaSalle, las trayectorias convergen cuando $t \to +\infty$ hacia el conjunto invariante maximal contenido en $\{ (u, v) \in \mathbb{R}^2 : v = 0 \}$, que son exactamente los puntos de equilibrio estático del péndulo:
    \[
        (u, v) = (k\pi, 0), \qquad k \in \mathbb{Z}.
    \]
    Para condiciones iniciales genéricas fuera de las variedades estables de los equilibrios inestables invertidos ($k$ impar), las oscilaciones se amortiguan y la solución converge asintóticamente a una posición de reposo colgado hacia abajo ($u(t) \to 2k\pi$, $u'(t) \to 0$).
}

% Ejercicio 2
\ejercicio{resuelto}{}{
    Demuéstrese que, para cada $p > 1$ y $x > 0$, el problema de valores iniciales
    \[
        \begin{cases}
            u'' = |u|^p, \\
            u(0) = x, \\
            u'(0) = 0,
        \end{cases}
    \]
    posee una única solución no prolongable estrictamente a la derecha que explota en tiempo finito, y calcúlese su tiempo de explosión analizando su dependencia con respecto de los parámetros $x > 0$ y $p > 1$.
}{
    Llamamos $(P)$ al problema de valor inicial.

    \textbf{1. Formulación como sistema plano y carácter Lipschitziano local:}

    Introduciendo la velocidad $v = u'$, formulamos $(P)$ como el sistema plano autónomo:
    \[
        \begin{cases}
            u' = v, \\
            v' = |u|^p,
        \end{cases}
        \iff
        \begin{pmatrix} u \\ v \end{pmatrix}' = \begin{pmatrix} v \\ |u|^p \end{pmatrix} = f(u, v), \qquad
        \begin{pmatrix} u(0) \\ v(0) \end{pmatrix} = \begin{pmatrix} x \\ 0 \end{pmatrix}.
    \]
    Dado que $p > 1$, la función $u \mapsto |u|^p = (u^2)^{p/2}$ es de clase $\mathcal{C}^1(\mathbb{R})$ con derivada continua:
    \[
        \frac{d}{du} |u|^p = p |u|^{p-1} \operatorname{sgn}(u).
    \]
    Por consiguiente, el campo $f \in \mathcal{C}^1(\mathbb{R}^2; \mathbb{R}^2)$ tiene matriz jacobiana continua:
    \[
        D_{(u,v)} f(u, v) = \begin{pmatrix} 0 & 1 \\ p |u|^{p-1} \operatorname{sgn}(u) & 0 \end{pmatrix}.
    \]
    Esto implica que $f$ es \textbf{localmente Lipschitziano} en $\mathbb{R}^2$, uniformemente respecto a $t \in \mathbb{R}$. Como $(x, 0) \in \R^2 = \Omega$, aplicamos los teoremas en cascada:
    \begin{enumerate}[label=\Roman*)]
        \item Por el Teorema de Cauchy-Lipschitz Local, existe una única solución local con dato inicial $(u(0), v(0)) = (x, 0)$.
        \item Por el Teorema de Unicidad Global, la solución es única en cualquier intervalo donde esté definida.
        \item Por el Teorema de Existencia de Solución Maximal, existe una \textbf{única solución maximal} $(I, u, v) = (I, u, u')$ no prolongable estrictamente a la derecha, con intervalo hacia el futuro:
    \end{enumerate}
    \[
        I = [0, T_{\max}), \qquad T_{\max} \le +\infty.
    \]
    Por el Segundo Teorema Fundamental, si $T_{\max} < +\infty$, entonces $\lim_{t \nearrow T_{\max}} \bigl( |u(t)| + |u'(t)| \bigr) = +\infty$.

    \medskip

    \textbf{2. Monotonía, positividad y convexidad:}

    Como $u(0) = x > 0$, por continuidad existe un primer intervalo donde $u(t) > 0$, de modo que $|u(t)|^p = u^p(t)$. Integrando la ecuación diferencial sobre $[0, t]$:
    \[
        u'(t) - \underbrace{u'(0)}_{= 0} = \int_0^t |u(s)|^p \, ds \ge 0 \implies u'(t) > 0 \qquad \forall t \in (0, T_{\max}).
    \]
    En consecuencia:
    \begin{itemize}
        \item $u(t)$ es \textbf{estrictamente creciente}, por lo que $u(t) > u(0) = x > 0$ para todo $t \in (0, T_{\max})$ (lo que confirma que $u(t)$ nunca se anula y la solución permanece positiva en todo $I$).
        \item $u''(t) = u^p(t) \ge x^p > 0$, luego $u(t)$ es \textbf{estrictamente convexa}.
        \item Como $u''(t) = v'(t) > 0$, la velocidad $v(t) = u'(t)$ es también \textbf{estrictamente creciente}.
    \end{itemize}
    Al ser ambas funciones monótonas crecientes y no negativas, la explosión de la norma implica que ambas variables divergen simultáneamente hacia $+\infty$:
    \[
        T_{\max} < +\infty \implies \lim_{t \nearrow T_{\max}} u(t) = +\infty \qquad \text{y} \qquad \lim_{t \nearrow T_{\max}} u'(t) = +\infty.
    \]

    \begin{center}
    \begin{tikzpicture}[>=stealth, scale=1.1]
        \draw[->] (-0.3,0) -- (3.8,0) node[right] {$t$};
        \draw[->] (0,-0.3) -- (0,3.6) node[above] {$u(t)$};
        
        \draw[dashed, gray] (2.8,0) node[below, black] {$T_{\max}$} -- (2.8,3.3);
        \draw[thick, blue!80!black, domain=0:2.55, samples=70]
            plot (\x, {0.8 + 0.1*\x + 0.35/((2.8 - \x)^1.2) - 0.35/(2.8^1.2)});
        
        \node[left] at (0,0.8) {$x$};
        \fill (0,0.8) circle (1.5pt);
        \node[font=\footnotesize, blue!80!black] at (1.4,2.5) {$u''(t) = u^p(t)$};
    \end{tikzpicture}
    \end{center}

    \textbf{3. Ley de conservación de energía y separación de variables:}

    Multiplicando la ecuación $u'' = u^p$ por la velocidad $u'$:
    \[
        u'(t) u''(t) = u^p(t) u'(t) \implies \frac{d}{dt} \left( \frac{v^2(t)}{2} - \frac{u^{p+1}(t)}{p+1} \right) = 0 \qquad \forall t \in I.
    \]
    Por consiguiente, la energía es constante e igual a su valor en $t = 0$:
    \[
        \frac{v^2(t)}{2} - \frac{u^{p+1}(t)}{p+1} = \frac{v^2(0)}{2} - \frac{u^{p+1}(0)}{p+1} = - \frac{x^{p+1}}{p+1} \qquad \forall t \in [0, T_{\max}).
    \]
    Despejando $v(t) = u'(t) > 0$:
    \[
        v^2(t) = \frac{2}{p+1} \left( u^{p+1}(t) - x^{p+1} \right) \implies v(t) = \sqrt{\frac{2}{p+1} \left( u^{p+1}(t) - x^{p+1} \right)}.
    \]
    Separando variables e integrando mediante el cambio $\zeta = u(s)$ (con $d\zeta = u'(s) ds$):
    \[
        t = \int_0^t 1 \, ds = \int_0^t \frac{u'(s)}{\sqrt{\frac{2}{p+1} \left( u^{p+1}(s) - x^{p+1} \right)}} \, ds = \int_x^{u(t)} \frac{d\zeta}{\sqrt{\frac{2}{p+1} \left( \zeta^{p+1} - x^{p+1} \right)}}.
    \]
    Efectuamos el cambio de escala $\zeta = x \theta$ (de modo que $d\zeta = x \, d\theta$ y los límites son $\theta = 1$ y $\theta = \frac{u(t)}{x}$):
    \[
        t = \int_1^{u(t)/x} \frac{x \, d\theta}{\sqrt{\frac{2}{p+1} x^{p+1} (\theta^{p+1} - 1)}} = \sqrt{\frac{p+1}{2}} x^{1 - \frac{p+1}{2}} \int_1^{u(t)/x} \frac{d\theta}{\sqrt{\theta^{p+1} - 1}} = \sqrt{\frac{p+1}{2}} x^{-\frac{p-1}{2}} \int_1^{u(t)/x} \frac{d\theta}{\sqrt{\theta^{p+1} - 1}}.
    \]
    Como la función integrando es continua y estrictamente positiva en $(1, +\infty)$, la integral define una función estrictamente creciente de $\frac{u(t)}{x}$, luego invertible.

    \medskip

    \textbf{4. Finitud y cálculo del tiempo de explosión $T_{\max}$:}

    Veamos primero que $u(t) \to +\infty$ cuando $t \nearrow T_{\max}$, sin suponer de antemano que $T_{\max} < +\infty$:
    \begin{itemize}
        \item Si $T_{\max} < +\infty$, lo hemos visto en el apartado 2 (Segundo Teorema Fundamental más monotonía).
        \item Si $T_{\max} = +\infty$, como $u''(t) = u^p(t) \ge x^p$ para todo $t \ge 0$, integrando dos veces y usando $u'(0) = 0$, $u(0) = x$:
        \[
            u'(t) \ge x^p t, \qquad u(t) \ge x + x^p \frac{t^2}{2} \xrightarrow{t \to +\infty} +\infty.
        \]
    \end{itemize}
    En ambos casos $\lim_{t \nearrow T_{\max}} u(t) = +\infty$. Tomando el límite cuando $t \nearrow T_{\max}$ en la expresión de $t$, el límite superior de integración tiende a $+\infty$:
    \[
        T_{\max} = \sqrt{\frac{p+1}{2}} \, x^{-\frac{p-1}{2}} \int_1^\infty \frac{d\theta}{\sqrt{\theta^{p+1} - 1}}.
    \]
    Si la integral impropia converge, el lado derecho es finito y, por tanto, $T_{\max} < +\infty$ (si fuese $T_{\max} = +\infty$ tendríamos $+\infty$ igual a un número finito, lo cual es absurdo). Demostremos que la integral impropia converge. Descomponemos para cualquier $\varepsilon > 0$ suficientemente pequeño y $M > 1 + \varepsilon$:
    \[
        \int_1^\infty \frac{d\theta}{\sqrt{\theta^{p+1} - 1}} = \int_1^{1+\varepsilon} \frac{d\theta}{\sqrt{\theta^{p+1} - 1}} + \int_{1+\varepsilon}^M \frac{d\theta}{\sqrt{\theta^{p+1} - 1}} + \int_M^\infty \frac{d\theta}{\sqrt{\theta^{p+1} - 1}}.
    \]
    \begin{itemize}
        \item \textit{Convergencia en el extremo inferior $\theta = 1$:} Consideremos la función $g(\theta) = \theta^{p+1} - 1$. Por el desarrollo de Taylor en $\theta = 1$:
        \[
            g(1) = 0, \qquad g'(1) = p + 1 > 0 \implies g(\theta) = (p + 1)(\theta - 1) + h(\theta)(\theta - 1) = \bigl( p + 1 + h(\theta) \bigr)(\theta - 1),
        \]
        donde $\lim_{\theta \to 1} h(\theta) = 0$. Tomando $\varepsilon > 0$ suficientemente pequeño, se tiene $p + 1 + h(\theta) \ge \frac{p+1}{2}$ para todo $\theta \in [1, 1+\varepsilon]$. Por tanto:
        \[
            \int_1^{1+\varepsilon} \frac{d\theta}{\sqrt{\theta^{p+1} - 1}} \le \sqrt{\frac{2}{p+1}} \int_1^{1+\varepsilon} \frac{d\theta}{\sqrt{\theta - 1}} = \sqrt{\frac{2}{p+1}} \Bigl[ 2\sqrt{\theta - 1} \Bigr]_1^{1+\varepsilon} = 2 \sqrt{\frac{2\varepsilon}{p+1}} < +\infty.
        \]
        \item \textit{Tramo intermedio $[1+\varepsilon, M]$:} La función integrando es continua sobre un compacto, luego la integral es finita.
        \item \textit{Convergencia en $+\infty$:} Para $\theta \to \infty$, el radicando se comporta como $\theta^{\frac{p+1}{2}}$. Como $p > 1$:
        \[
            \frac{p+1}{2} > \frac{1+1}{2} = 1.
        \]
        Fijando $M > 1$ tal que $\theta^{p+1} - 1 \ge \frac{1}{2} \theta^{p+1}$ para todo $\theta \ge M$:
        \[
            \int_M^\infty \frac{d\theta}{\sqrt{\theta^{p+1} - 1}} \le \sqrt{2} \int_M^\infty \theta^{-\frac{p+1}{2}} \, d\theta = \sqrt{2} \left[ \frac{\theta^{-\frac{p-1}{2}}}{-\frac{p-1}{2}} \right]_M^\infty = \frac{2\sqrt{2}}{(p-1) M^{\frac{p-1}{2}}} < +\infty.
        \]
    \end{itemize}
    Queda demostrado que la integral impropia converge a una constante positiva:
    \[
        J(p) = \int_1^\infty \frac{d\theta}{\sqrt{\theta^{p+1} - 1}} < +\infty.
    \]
    En consecuencia, la solución explota en tiempo finito con:
    \[
        T_{\max}(x, p) = \sqrt{\frac{p+1}{2}} \, J(p) \, x^{-\frac{p-1}{2}} < +\infty.
    \]

    Nótese que se podría haber aplicado el teorema del valor medio a la función:
    \[
        f(\theta) = \theta^{p+1} \implies
        \theta^{p+1} - 1 = f(\theta) - f(1) = f'(\xi)(\theta - 1) = (p + 1) \xi^p (\theta - 1) \ge (p + 1) (\theta - 1), \qquad \xi \in (1, \theta),
    \]
    ya que $\xi > 1$ implica $\xi^p > 1$.
    Permitiendonos de nuevo acotar la integral por la de $(\theta - 1)^{-1/2}$ y concluir la convergencia.

    \medskip

    \textbf{5. Dependencia respecto de los parámetros $x > 0$ y $p > 1$:}
    \begin{itemize}
        \item \textbf{Dependencia respecto de $x$:} Como $-\frac{p-1}{2} < 0$, la función $x \mapsto T_{\max}(x, p)$ es estrictamente decreciente:
        \[
            \lim_{x \to 0^+} T_{\max}(x, p) = +\infty, \qquad \lim_{x \to +\infty} T_{\max}(x, p) = 0.
        \]
        A mayor dato inicial $x$, la explosión se produce de manera mucho más rápida.
        \item \textbf{Dependencia respecto de $p$:} Cuando $p \to 1^+$, el exponente de decaimiento $-\frac{p-1}{2} \to 0$ y la integral en el infinito $\int_M^\infty \theta^{-\frac{p+1}{2}} d\theta$ diverge logarítmicamente (puesto que para la ecuación lineal $u'' = u$, las soluciones son combinaciones de exponenciales $e^{\pm t}$ globales sin explosión en tiempo finito). Por consiguiente:
        \[
            \lim_{p \to 1^+} T_{\max}(x, p) = +\infty.
        \]
        Para potencias $p$ grandes, la no linealidad acelera intensamente la tasa de divergencia temporal.
    \end{itemize}

    \textbf{Justificación de los límites respecto de $p$.} Usando el Teorema del Valor Medio anterior y $\theta^{p+1} - 1 \le \theta^{p+1}$, obtenemos cotas explícitas de $J(p)$:
    \begin{itemize}
        \item \textit{Cota inferior:} como $\theta^{p+1} - 1 \le \theta^{p+1}$,
        \[
            J(p) \ge \int_1^\infty \theta^{-\frac{p+1}{2}} \, d\theta = \left[ \frac{\theta^{-\frac{p-1}{2}}}{-\frac{p-1}{2}} \right]_1^\infty = \frac{2}{p-1}.
        \]
        Por tanto
        \[
            T_{\max}(x, p) \ge \sqrt{\frac{p+1}{2}} \, \frac{2}{p-1} \, x^{-\frac{p-1}{2}} \xrightarrow{p \to 1^+} +\infty,
        \]
        pues $\sqrt{\frac{p+1}{2}} \to 1$, $x^{-\frac{p-1}{2}} \to 1$ y $\frac{2}{p-1} \to +\infty$. Esto demuestra rigurosamente que $\lim_{p \to 1^+} T_{\max}(x, p) = +\infty$.
        \item \textit{Cota superior:} en $[1, 2]$ usamos $\theta^{p+1} - 1 \ge (p+1)(\theta - 1)$, y en $[2, +\infty)$ usamos $\theta^{p+1} - 1 \ge \frac{1}{2} \theta^{p+1}$ (cierto porque $\theta^{p+1} \ge 2^{p+1} \ge 2$):
        \[
            J(p) \le \frac{1}{\sqrt{p+1}} \int_1^2 \frac{d\theta}{\sqrt{\theta - 1}} + \sqrt{2} \int_2^\infty \theta^{-\frac{p+1}{2}} \, d\theta = \frac{2}{\sqrt{p+1}} + \frac{2\sqrt{2}}{p-1} \, 2^{-\frac{p-1}{2}}.
        \]
        Así, $\sqrt{\frac{p+1}{2}} \, J(p) \le \sqrt{2} + \frac{2\sqrt{p+1}}{p-1} \, 2^{-\frac{p-1}{2}}$, que está acotado para $p \ge 2$.
    \end{itemize}
    \textbf{Observación (cuidado con $p \to +\infty$):} la afirmación de que $p$ grande acelera la explosión solo es cierta si $x > 1$. En efecto, el factor $\sqrt{\frac{p+1}{2}} \, J(p)$ está acotado, así que todo depende de $x^{-\frac{p-1}{2}}$:
    \begin{itemize}
        \item Si $x > 1$: $x^{-\frac{p-1}{2}} \to 0$ exponencialmente, luego $T_{\max}(x, p) \xrightarrow{p \to \infty} 0$.
        \item Si $0 < x < 1$: por la cota inferior, $T_{\max}(x, p) \ge \sqrt{\frac{p+1}{2}} \frac{2}{p-1} \, x^{-\frac{p-1}{2}}$, y la exponencial $x^{-\frac{p-1}{2}} \to +\infty$ domina al factor $\frac{\sqrt{p+1}}{p-1} \sim \frac{1}{\sqrt{p}}$, luego $T_{\max}(x, p) \xrightarrow{p \to \infty} +\infty$. Intuitivamente, si $0 < u < 1$, entonces $u^p$ es muy pequeño para $p$ grande y la solución tarda mucho en despegar.
    \end{itemize}
}

% Ejercicio 3
\ejercicio{medio}{}{
    Dados $\lambda > 0$, $m > 0$, $p > 1$ y $x \ge 0$, demuéstrese que el problema
    \[
        \begin{cases}
            -u'' = \lambda u - m|u|^p, \\
            u(0) = x, \\
            u'(0) = 0,
        \end{cases}
    \]
    posee una única solución maximal a la derecha, y estúdiese su naturaleza de acuerdo con los diferentes valores que toma $x$.
}{
    Llamamos $(P)$ al problema de valor inicial.

    \textbf{1. Formulación como sistema de primer orden y carácter Lipschitziano local:}

    Introduciendo la velocidad $v = u'$, el problema $(P)$ equivale al sistema de primer orden asociado
    \[
        \begin{cases}
            u' = v, \\
            v' = m|u|^p - \lambda u, \\
            (u(0), v(0)) = (x, 0),
        \end{cases}
        \qquad
        f(u, v) = \begin{pmatrix} v \\ m|u|^p - \lambda u \end{pmatrix} \in \mathcal{C}^1(\R^2; \R^2).
    \]
    Su matriz jacobiana es
    \[
        D_{(u,v)} f(u, v) = \begin{pmatrix} 0 & 1 \\ p m |u|^{p-1} \operatorname{signo}(u) - \lambda & 0 \end{pmatrix} \in \mathcal{C}(\R^2; \mathcal{M}_2(\R)),
    \]
    que es continua en todo $\R^2$: el único punto delicado es $u = 0$, pero como $p > 1$ se tiene
    \[
        |u|^{p-1} \xrightarrow{u \to 0} 0,
    \]
    luego $p m |u|^{p-1} \operatorname{signo}(u) \to 0$ y la entrada es continua también en $u = 0$. Por tanto, $f$ es localmente Lipschitziana respecto de $(u, v) \in \R^2$, uniformemente en $t \in \R$ (el sistema es autónomo).

    \textbf{2. Existencia y unicidad de la solución maximal:}

    Como $(x, 0) \in \R^2 = \Omega$, aplicamos los teoremas en cascada:
    \begin{enumerate}[label=\Roman*)]
        \item Por el Teorema de Cauchy-Lipschitz Local, existe una única solución local con dato inicial $(u(0), v(0)) = (x, 0)$.
        \item Por el Teorema de Unicidad Global, la solución es única en cualquier intervalo donde esté definida.
        \item Por el Teorema de Existencia de Solución Maximal, existe una única solución maximal a la derecha $(I, u, v)$.
    \end{enumerate}
    Por el Segundo Teorema Fundamental (pues el campo está definido en todo $\R^2$), $I = [0, T_{\max})$ con $T_{\max} \le +\infty$, y
    \[
        T_{\max} < +\infty \implies \lim_{t \nearrow T_{\max}} \big( |u(t)| + |v(t)| \big) = +\infty.
    \]

    \textbf{3. Ley de conservación de la energía:}

    Para $t \in I$, multiplicamos la ecuación $v' = m|u|^p - \lambda u$ por $v = u'$:
    \[
        \big( m|u(t)|^p - \lambda u(t) \big) u'(t) = v'(t) v(t).
    \]
    Ambos miembros son derivadas exactas. Usando que $\frac{d}{dz} \left( |z|^{p+1} \operatorname{signo}(z) \right) = (p+1)|z|^p$:
    \[
        \left( \frac{v^2(t)}{2} + \frac{\lambda u^2(t)}{2} - \frac{m |u(t)|^{p+1} \operatorname{signo}(u(t))}{p+1} \right)' = 0 \quad \forall t \in I.
    \]
    Por tanto, la expresión es constante e igual a su valor en $t = 0$, donde $u(0) = x \ge 0$ y $v(0) = 0$:
    \[
        \frac{v^2(t)}{2} + \frac{\lambda u^2(t)}{2} - \frac{m}{p+1} \operatorname{signo}(u(t)) |u(t)|^{p+1} = \frac{\lambda x^2}{2} - \frac{m}{p+1} x^{p+1}.
    \]
    Definimos el \textbf{potencial}: con el signo menos de la ecuación $-u'' = \lambda u - m|u|^p$, el potencial es la primitiva del miembro derecho,
    \[
        \varphi(z) = \frac{\lambda z^2}{2} - \frac{m}{p+1} \operatorname{signo}(z) |z|^{p+1} =
        \begin{cases}
            \dfrac{\lambda z^2}{2} - \dfrac{m}{p+1} z^{p+1} & \text{si } z \ge 0, \\[1em]
            \dfrac{\lambda z^2}{2} + \dfrac{m}{p+1} |z|^{p+1} & \text{si } z < 0,
        \end{cases}
        \qquad \varphi'(z) = \lambda z - m |z|^p.
    \]
    Con esta notación, la ley de conservación se escribe
    \[
        \boxed{\frac{v^2(t)}{2} + \varphi(u(t)) = \varphi(x) =: E \quad \forall t \in I.}
    \]

    \textbf{4. Estudio del potencial $\varphi$:}
    \begin{itemize}
        \item \textit{Signo de $\varphi'$ para $z < 0$:} $\varphi'(z) = \lambda z - m|z|^p < 0$, pues $\lambda z < 0$ y $-m|z|^p < 0$. Así, $\varphi$ es estrictamente decreciente en $(-\infty, 0)$, con $\varphi(z) \to +\infty$ cuando $z \to -\infty$.

        \item \textit{Ceros de $\varphi$ en $z \ge 0$:}
        \[
            \frac{\lambda z^2}{2} - \frac{m}{p+1} z^{p+1} = 0 \iff
            \begin{cases}
                z = 0, \\
                \dfrac{\lambda}{2} - \dfrac{m}{p+1} z^{p-1} = 0 \implies z = \left( \dfrac{\lambda (p+1)}{2m} \right)^{\frac{1}{p-1}} =: \alpha.
            \end{cases}
        \]

        \item \textit{Puntos críticos de $\varphi$ en $z \ge 0$:}
        \[
            \lambda z - m z^p = 0 \iff
            \begin{cases}
                z = 0, \\
                z^{p-1} = \dfrac{\lambda}{m} \implies z = \left( \dfrac{\lambda}{m} \right)^{\frac{1}{p-1}} =: \omega.
            \end{cases}
        \]
        Como $p > 1$ se tiene $\frac{p+1}{2} > 1$, luego $\frac{\lambda(p+1)}{2m} > \frac{\lambda}{m}$ y, por tanto, $\omega < \alpha$. Además, para $z > 0$, $\varphi'(z) = z(\lambda - m z^{p-1})$ es positiva en $(0, \omega)$ y negativa en $(\omega, +\infty)$.

        \item \textit{Naturaleza de los puntos críticos:} para $z \ge 0$, $\varphi''(z) = \lambda - m p z^{p-1}$, luego
        \[
            \varphi''(0) = \lambda > 0 \quad (\text{mínimo local}), \qquad
            \varphi''(\omega) = \lambda - m p \frac{\lambda}{m} = (1 - p) \lambda < 0 \quad (\text{máximo local}),
        \]
        ya que $1 - p < 0$ y $\lambda > 0$.
    \end{itemize}

    \begin{center}
    \begin{tikzpicture}[>=stealth, xscale=2.6, yscale=5.5]
        % ejes
        \draw[->] (-1.0,0) -- (1.95,0) node[right] {$z$};
        \draw[->] (0,-0.3) -- (0,0.72) node[above] {$\varphi(z)$};
        % potencial (lambda = m = 1, p = 2)
        \draw[thick, blue!80!black, domain=-0.95:0, samples=60] plot (\x, {\x*\x/2 - \x*\x*\x/3});
        \draw[thick, blue!80!black, domain=0:1.72, samples=80] plot (\x, {\x*\x/2 - \x*\x*\x/3});
        % nivel phi(x)
        \draw[dashed, gray] (-1.0,0.108) node[left, black] {$\varphi(x)$} -- (1.9,0.108);
        % puntos
        \draw[dotted] (-0.4117,0.108) -- (-0.4117,0) node[below] {$\beta(x)$};
        \draw[dotted] (0.6,0.108) -- (0.6,0) node[below] {$x$};
        \draw[dotted] (1,0.1667) -- (1,0) node[below] {$\omega$};
        \draw[dotted] (1.3117,0.108) -- (1.3117,0) node[below] {$x$};
        \fill (1.5,0) circle (0.25pt) node[below right] {$\alpha$};
        % conjuntos donde puede moverse u
        \draw[very thick, red!80!black] (-0.4117,0) -- (0.6,0);
        \draw[very thick, red!80!black] (1.3117,0) -- (1.9,0);
    \end{tikzpicture}
    \end{center}
    En rojo, los intervalos donde puede moverse $u(t)$: $[\beta(x), x]$ si $x \in (0, \omega)$, y $[x, +\infty)$ si $x > \omega$. Aquí $\beta(x) < 0$ es el único punto negativo con $\varphi(\beta(x)) = \varphi(x)$.

    \textbf{5. Consecuencia de la ley de conservación:}

    Como $\frac{v^2(t)}{2} \ge 0$,
    \[
        \varphi(u(t)) = \varphi(x) - \frac{v^2(t)}{2} \le \varphi(x) \quad \forall t \in I.
    \]
    Es decir, la solución solo puede moverse por la región $\{ z \in \R : \varphi(z) \le \varphi(x) \}$, y, al ser $u$ continua con $u(0) = x$, por conexión se queda en la componente de ese conjunto que contiene a $x$.

    \textbf{6. Casos según el dato inicial $x$:}
    \begin{itemize}
        \item \textbf{Si $x = 0$:} $(u, v) \equiv (0, 0)$ es solución, pues $f(0, 0) = (0, 0)$. Por unicidad, $u \equiv 0$.

        \item \textbf{Si $x = \omega$:} $f(\omega, 0) = (0, m\omega^p - \lambda \omega) = (0, \omega(m\omega^{p-1} - \lambda)) = (0, 0)$, así que $(\omega, 0)$ es un equilibrio, y por unicidad $u \equiv \omega$.

        \item \textbf{Si $x \in (0, \omega)$:} el conjunto $\{ \varphi \le \varphi(x) \}$ cerca de $x$ es $[\beta(x), x]$, separado del resto del conjunto porque $\varphi(\omega) > \varphi(x)$. Luego $u(t) \in [\beta(x), x]$ está acotada. Entonces $v(t)$ también lo está, ya que $\varphi \ge 0$ en $[\beta(x), x]$ (su mínimo es $\varphi(0) = 0$):
        \[
            v^2(t) = 2 \big( \varphi(x) - \varphi(u(t)) \big) \le 2 \varphi(x) \implies |v(t)| \le \sqrt{2 \varphi(x)} \quad \forall t \in I.
        \]
        Así, $|u(t)| + |v(t)|$ está acotada en $I$ y, por el Segundo Teorema Fundamental, $T_{\max} = +\infty$: la solución está globalmente definida.

        \item \textbf{Si $x \in (\omega, +\infty)$:} como $\varphi$ es decreciente en $(\omega, +\infty)$ y $\varphi(\omega) > \varphi(x)$, de $\varphi(u(t)) \le \varphi(x)$ se deduce
        \[
            u(t) \ge x \quad \forall t \in I.
        \]
        En particular $u(t) > 0$, así que $|u(t)|^p = u^p(t)$ y la ecuación queda
        \[
            -u''(t) = \big( \lambda - m u^{p-1}(t) \big) u(t) \overset{(\ast)}{<} 0.
        \]
        $(\ast)$: como $u(t) \ge x > \omega$,
        \[
            u(t) > \left( \frac{\lambda}{m} \right)^{\frac{1}{p-1}} \implies u^{p-1}(t) > \frac{\lambda}{m} \implies \lambda - m u^{p-1}(t) < 0,
        \]
        y además $u(t) > 0$. Por tanto $u''(t) > 0$, así que $u'(t)$ es estrictamente creciente y, como $u'(0) = 0$,
        \[
            u'(t) > 0 \quad \forall t \in (0, T_{\max}).
        \]

            Puesto que $u(t) \ge x > \omega$, de la ley de conservación deducimos:
        \[
            \frac{v^2(t)}{2} + \varphi(u(t)) = \varphi(x) \implies v^2(t) = 2\big(\varphi(x) - \varphi(u(t))\big).
        \]
        Sustituyendo la expresión del potencial $\varphi(z) = \frac{\lambda}{2} z^2 - \frac{m}{p+1} z^{p+1}$:
        \[
            v(t) = \sqrt{\frac{2m}{p+1}\big(u^{p+1}(t) - x^{p+1}\big) - \lambda\big(u^2(t) - x^2\big)}.
        \]
        Como $v(t) = u'(t) > 0$, integramos por separación de variables en $[0, t]$ realizando el cambio $z = u(s)$ y posteriormente el reescalado $z = x\theta$ ($dz = x\,d\theta$):
        \[
            t = \int_0^t dt = \int_x^{u(t)} \frac{dz}{\sqrt{\frac{2m}{p+1}(z^{p+1} - x^{p+1}) - \lambda(z^2 - x^2)}} = \int_1^{u(t)/x} \frac{d\theta}{\sqrt{\frac{2m x^{p-1}}{p+1}(\theta^{p+1} - 1) - \lambda(\theta^2 - 1)}}.
        \]
        Definiendo el radicando $f(\theta) = \frac{2m x^{p-1}}{p+1}(\theta^{p+1} - 1) - \lambda(\theta^2 - 1)$, el tiempo maximal de existencia hacia el futuro está dado por el límite cuando $u(t) \to +\infty$:
        \[
            T_{\max}(x, p) = \int_1^{+\infty} \frac{d\theta}{\sqrt{f(\theta)}} = \int_1^M \frac{d\theta}{\sqrt{f(\theta)}} + \int_M^{+\infty} \frac{d\theta}{\sqrt{f(\theta)}} \qquad (M > 1).
        \]

        \textbf{Convergencia de la integral en $\theta \to 1^+$:}
        
        Observamos que $f(1) = 0$. Derivando respecto de $\theta$:
        \[
            f'(\theta) = 2m x^{p-1}\theta^p - 2\lambda\theta = 2\theta(m x^{p-1}\theta^{p-1} - \lambda).
        \]
        Para $\theta > 1$ y siendo $x > \omega = \left(\frac{\lambda}{m}\right)^{\frac{1}{p-1}}$, tenemos $m x^{p-1} - \lambda > 0$, luego $f'(\theta) > 2\theta(m x^{p-1} - \lambda) > 0$. Además:
        \[
            f''(\theta) = 2m p x^{p-1}\theta^{p-1} - 2\lambda > 2(m x^{p-1} - \lambda) > 2\left(m \frac{\lambda}{m} - \lambda\right) = 0,
        \]
        por lo que $f'(\theta)$ es estrictamente creciente para $\theta \ge 1$. Por el desarrollo de Taylor con resto de Lagrange (o el Teorema del Valor Medio), existe $\xi \in (1, \theta)$ tal que:
        \[
            f(\theta) = f(1) + f'(\xi)(\theta - 1) \ge f'(1)(\theta - 1) = 2(mx^{p-1} - \lambda)(\theta - 1).
        \]
        Por consiguiente:
        \[
            \int_1^M \frac{d\theta}{\sqrt{f(\theta)}} \le \frac{1}{\sqrt{2(mx^{p-1} - \lambda)}} \int_1^M (\theta - 1)^{-1/2}\,d\theta = \frac{2\sqrt{M-1}}{\sqrt{2(mx^{p-1} - \lambda)}} < +\infty.
        \]

        \textbf{Convergencia de la integral en $\theta \to +\infty$:}
        
        Para la cola asintótica cuando $\theta \to +\infty$, el término dominante del radicando es de orden $\theta^{p+1}$:
        \[
            \lim_{\theta \to +\infty} \frac{f(\theta)}{\frac{2m x^{p-1}}{p+1}\theta^{p+1}} = \lim_{\theta \to +\infty} \left(1 - \frac{1}{\theta^{p+1}} - \frac{\lambda(p+1)}{2mx^{p-1}}\left(\frac{1}{\theta^{p-1}} - \frac{1}{\theta^{p+1}}\right)\right) = 1.
        \]
        Por tanto, fijando $M$ suficientemente grande, $\sqrt{f(\theta)} \ge \frac{1}{2}\sqrt{\frac{2m x^{p-1}}{p+1}}\,\theta^{\frac{p+1}{2}}$, de donde:
        \[
            \int_M^{+\infty} \frac{d\theta}{\sqrt{f(\theta)}} \le 2\sqrt{\frac{p+1}{2mx^{p-1}}} \int_M^{+\infty} \theta^{-\frac{p+1}{2}}\,d\theta.
        \]
        Al ser $p > 1$, se cumple $\frac{p+1}{2} > 1$, luego el exponente $-\frac{p+1}{2} < -1$ y la integral de Riemann converge:
        \[
            \int_M^{+\infty} \theta^{-\frac{p+1}{2}}\,d\theta = \left[\frac{\theta^{1 - \frac{p+1}{2}}}{1 - \frac{p+1}{2}}\right]_M^{+\infty} = \frac{M^{-\frac{p-1}{2}}}{\frac{p-1}{2}} < +\infty.
        \]
        Concluimos que $T_{\max}(x, p) < +\infty$, por lo que la solución \textbf{explota en tiempo finito hacia $+\infty$}:
        \[
            \lim_{t \to T_{\max}^-} u(t) = +\infty, \qquad \lim_{t \to T_{\max}^-} u'(t) = +\infty.
        \]
        Además, analizando la dependencia de $T_{\max}(x, p)$ respecto de $x$, cuando $x \to \omega^+ = \left(\frac{\lambda}{m}\right)^{\frac{1}{p-1}}$, el denominador en $\theta = 1$ verifica $f'(1) = 2(mx^{p-1} - \lambda) \to 0$, por lo que la singularidad en $\theta = 1$ se vuelve no integrable (cero de orden doble con comportamiento local asintótico como $\frac{1}{\theta-1}$), obteniéndose $\lim_{x \to \omega^+} T_{\max}(x, p) = +\infty$, consistente con que $u \equiv \omega$ es una solución de equilibrio global. Por contra, cuando $x \to +\infty$, $T_{\max}(x, p) \to 0^+$.

        \begin{center}
        \begin{tikzpicture}[scale=1.1]
            \draw[->] (-0.3,0) -- (4.5,0) node[right] {$x$};
            \draw[->] (0,-0.3) -- (0,3.5) node[above] {$T_{\max}(x)$};
            \draw[dashed, blue!70!black] (1.2,0) -- (1.2,3.3);
            \node[below] at (1.2,0) {$\left(\frac{\lambda}{m}\right)^{\frac{1}{p-1}}$};
            \draw[thick, blue!80!black, domain=1.35:4.2, samples=100] plot (\x, {0.6/(\x - 1.2) + 0.15});
        \end{tikzpicture}
        \end{center}
    \end{itemize}
}

% Ejercicio 4
\ejercicio{medio}{}{
    Caracterizar las funciones $f \in \mathcal{C}^1([0, +\infty); \R)$ tales que $f(0) = 0$ y $f'(u) > 0$ para cada $u > 0$ para las que la solución del problema de Cauchy
    \[
        \begin{cases}
            u'' = f(u), \\
            u(0) = x > 0, \\
            u'(0) = 0,
        \end{cases}
    \]
    explota en tiempo finito.
    
    Dado $R > 0$, analizar la existencia de soluciones positivas para el problema singular de valores de contorno
    \[
        \begin{cases}
            u'' = f(u) \quad \text{en } (-R, R), \\
            u(-R) = +\infty = u(R).
        \end{cases}
    \]
}{
    \textbf{1. Formulación como sistema de primer orden y conservación de energía:}

    Introduciendo $v = u'$, el problema de Cauchy para la ecuación de segundo orden se formula como el sistema autónomo en el plano de fases:
    \[
        \begin{pmatrix} u \\ v \end{pmatrix}' = \begin{pmatrix} v \\ f(u) \end{pmatrix}, \qquad \begin{pmatrix} u(0) \\ v(0) \end{pmatrix} = \begin{pmatrix} x \\ 0 \end{pmatrix}.
    \]
    Como $f \in \mathcal{C}^1([0, +\infty); \R)$, el campo es localmente Lipschitziano en el semiplano $u \ge 0$. Por el Teorema Fundamental de existencia y unicidad maximal, existe una única solución maximal $(I, u)$ con $I = [0, T_{\max})$.

    Multiplicando la ecuación escalar por $u'(t) = v(t)$:
    \[
        u''(t) u'(t) = f(u(t)) u'(t) \implies \left( \frac{v^2(t)}{2} - \int_x^{u(t)} f(s)\,ds \right)' = 0 \quad \forall t \in I.
    \]

    \begin{center}
        \fcolorbox{orange!85!black}{orange!10}{\parbox{0.92\linewidth}{\centering
            \textbf{Hasta aquí llegamos en clase (10/11).}\\[0.2em]
            Se continúa completando la caracterización de explosión en tiempo finito de la integral singular en la próxima sesión.
        }}
    \end{center}
}

% Ejercicio 5
\ejercicio{enunciado}{}{
    Analizar la existencia, unicidad y prolongabilidad de las soluciones del problema de valores iniciales para el sistema de competición difusivo unidimensional
    \[
        \begin{cases}
            -u'' = \lambda u - a u^2 - b u v, \\
            -v'' = \mu v - d v^2 - c u v, \\
            u(0) = x, \quad u'(0) = 0, \\
            v(0) = y, \quad v'(0) = 0,
        \end{cases}
    \]
    donde $\lambda, \mu, a, b, c, d, x$ e $y$ son constantes positivas.
}{}
